% =============================================================================
%  Ouverture — Seraphina
%  De son arrivée chez les Whitmore, l'après-midi, jusqu'à la lecture de la lettre.
% =============================================================================
\partie{Ouverture}{Seraphina}

% Rappels d'écriture :
%   \separation            changement de scène
%   \areprendre{...}       passage à retravailler plus tard, sans s'arrêter
%   \noteauteur{...}       note pour soi, masquée en version finale
%   \begin{lettre} ... \end{lettre}, \begin{telegramme}..., \begin{coupure}...
%   « guillemets » : taper directement « et », ou \enquote{...}
%   Dialogues : \rep en début de chaque réplique (une par changement de locuteur).

\chapitre{Valse (provisoire)}
% Touver un nom qui evoque l'idée de l'opera, piece, valse, comedie, danse, qui evoque l'idée que toutes ces obligations sociiales ne sont qu'un sepcatcle de facade

\noteauteur{Réfléchir a la manière dont Lady Whitmore s'adresse a Sera. le tutoiement serait une marque d'intimité forte. Sera, elle, doit utiliser le vouvoiement.}

\debut{B}{elgravia}, Londres, 27 novembre 1890.
Alors que les neuf coups de la grande horloge du hall de la demeure des Whitmore résonnaient en cette froide soirée d'automne, Seraphina sentit monter en elle une angoisse qu'elle ne connaissait que trop bien.
Ces dernières semaines, les lettres signées de la main de la signora Valentini avaient connu un rythme inégalé depuis que cette Italienne s’était mise en tête de retrouver l’auteur du vol des dernières œuvres de feu son époux, sans doute en raison de son séjour à Londres.
En effet, la veuve semblait visiblement s'être mise en tête que Seraphina Blackwood, chanteuse reconnue dans tous les hauts cercles londoniens et membre du duo \enquote{Les Perles de Londres}, en était responsable, et comptait bien lui faire regretter cette action.

La simple pensée de cette vieille folle fit frissonner Seraphina.
Pourquoi diable fallut-il qu'elle soit présente à Londres un 27 novembre ?
Pourquoi fallait-il que cela tombe le jour de son anniversaire, alors que Lady Whitmore donnait une réception dont elle seule avait le secret en son honneur et en celui du duo qu'elle formait avec sa fille, Evangeline ?
Car le malheur de Seraphina fut que la présence d'une dame de ce rang à Londres, influente au sein des plus hauts cercles artistiques européens, était impossible à ignorer, et elle se vit contrainte d'accepter que Lady Whitmore l'invite si elle ne voulait pas passer pour l'inconvenante de service.
Et au grand malheur de Seraphina, Valentini avait accepté \enquote{avec une joie non dissimulée} de se joindre à la réception.
Elle fut à nouveau prise d'un profond frisson.
\rep Sera, tu as froid ?
Je te vois frissonner depuis quelques minutes, lança alors Lady Whitmore, l’hôtesse de cette soirée.
Je peux demander à un domestique de rajouter du charbon, reprit-elle.

Interrompue dans ses réflexions, Seraphina hésita quelques secondes avant de répondre.
\rep Non, tout va bien.
Simplement\ldots{} simplement le trac qui commence à monter\ldots{} je crois.
\rep Ne t'inquiète pas, ta prestation avec Eva sera parfaite, comme à votre habitude.

Alors qu'elle allait continuer, Lady Whitmore fut interrompue par le claquement sec et régulier des sabots d'un cheval sur les pavés des rues froides et sombres de Londres.
\rep Ah, je pense que nos premiers invités approchent, s'exclama-t-elle.
Je mettrais ma main à couper qu'il s'agit du Dr Pemberton et de son collègue photographe, ils sont réputés pour être d'une ponctualité à toute épreuve, reprit-elle.

Elle s'avança alors sur les premières marches de l'escalier.
\rep Eva, rejoins-nous dans le hall, je crois que nos premiers invités arrivent, lança-t-elle enfin, avant de réajuster rapidement sa coiffure devant la grande glace qu'elle avait fait venir de Venise.

\separation

Alors que la valse des invités s’enchaînait, Dr Pemberton et Theodore Bradshaw, Miss Beaumont, Miss Thornfield, Sir Edmund, Dr Ashford, Seraphina sentit que son corps lui échappait, comme si elle n’était plus qu'une spectatrice d'elle-même, comme si elle n’était plus que l'un de ces automates dont certains horlogers parisiens s’étaient faits spécialistes.
Plus le temps passait, plus les invités défilaient, plus Seraphina redoutait la prochaine annonce du majordome.

Tandis qu'elle quittait la salle à manger dans laquelle s’étaient réunis les premiers arrivés pour se rendre dans le hall et accueillir les Fairfax, qui venaient d'être annoncés, son regard se porta sur l'horloge du hall, cet implacable maître du temps, insensible aux futiles ressentis de l'âme humaine.
Vingt et une heures vingt-deux.
Vingt-deux minutes depuis vingt et une heures.
Cette épreuve qui semblait déjà avoir duré une bonne centaine d'éternités à Seraphina n'avait commencé qu'il y avait vingt-deux minutes.
Vingt-deux minutes de souffrance.
Vingt-deux minutes à s'imaginer les pires scénarios.
Vingt-deux minutes durant lesquelles chaque pas, chaque mot, chaque respiration lui demandait un effort surhumain.
Vingt-deux minutes à se demander si elle survivrait à la prochaine.
Et l'objet de sa souffrance n’était même pas encore présent en ces lieux.

Alors qu'elle terminait des salutations cordiales, mais sans chaleur, pour accueillir le couple Fairfax, Seraphina contint un sanglot.
Pourquoi avait-il fallu qu'elle dise oui à cet homme ce jour maudit ?
Pourquoi n'avait-elle pas pris plus de précautions avant de faiblement hocher la tête \noteauteur{S'assurer que cela n'en dit pas plus que souhaité sur les partitions} ?
Elle refréna de justesse un nouveau sanglot.
Il fallait à tout prix qu'elle maintienne les apparences, qu'elle ne laisse pas transparaître sa détresse.
Pour elle.
Et pour Eva.

Elle fit un discret signe de la main, et Florence, sa dame de compagnie, devenue avec le temps la grande sœur qu'elle n'avait jamais eue, vint immédiatement avec un verre de citronnade, et lui tendit délicatement un petit mouchoir de coton blanc brodé de ses initiales, \enquote{M.F.}.
Mary Fletcher.
Elle eut une pensée nostalgique pour la jeune fille qu'elle était avant, une simple orpheline de Whitechapel comme il y en avait tant d'autres.
Oh, la vie n'était pas simple dans ces rues crasseuses, bien au contraire !
Mais depuis quelque temps, Seraphina se surprenait à se demander ce qui se serait passé si elle n'avait pas franchi la porte du concours, ce jour-là.
Si elle n'avait pas rencontré Evangeline.
Si ensemble, elles n'avaient pas formé de duo.
Si elle était restée une anonyme qui n'avait qu'à se soucier de trouver suffisamment de nourriture pour survivre.
Elle chassa rapidement cette pensée de son esprit et prit le verre, mais refusa poliment le mouchoir.
Tandis qu'elle prenait prudemment une gorgée, le majordome annonça de sa voix claire et sans émotions :
\rep La signora Valentini !

Seraphina, prise de court lors de cet instant de douceur avec Florence, manqua de s’étrangler et faillit lâcher son verre de surprise.
\rep Que se passe-t-il ?
s'exclama-t-elle.
\rep La signora Valentini vient d'être annoncée, lui répondit Florence de sa voix douce.
Je pense que vous devriez vous joindre à Lady Whitmore pour l’accueillir, finit-elle.

\separation

\rep Seraphina ?
Seraphina ?
Vous m’écoutez ?
Que pensez-vous de cela ?

C’était Evangeline, aux côtés de Theodore, qui lui demandait son avis.
Sur quoi, elle n'en avait pas la moindre idée.

\rep Veuillez m'excuser, je\ldots{} je crains de n'avoir décroché quelques instants de cette discussion, répondit-elle.

Tout ce dont elle se souvenait, c’était l'annonce de l’arrivée de cette satanée Valentini.
\pensee{Que le diable la fasse rejoindre son époux}, s’exclama-t-elle en son for intérieur.

\rep Nous vous demandions si vous aviez une idée de ce qui pourrait pousser Maestro Eristoff à être absent ce soir, fit Theodore.
\noteauteur{Evangeline a les yeux rouges, et peut être une larme. Remarqué par Theodore, qui pourra le partager plus tard. TbD: Seraphina l'a t'elle remarqué malgré son état ?}

Alors qu'elle tournait la tête pour essayer de distinguer les différentes personnes présentes dans la salle à manger, décidément bien trop étouffante en ce mois de novembre, elle vit d'abord Lady Eristoff, dont elle n’arrivait pas à se remémorer l'arrivée, puis elle l’aperçut !
\rep Je\ldots{} je n'en ai pas la moindre idée, articula-t-elle avec difficulté.
\rep Sera, tu vas bien ? \noteauteur{Evangeline switch vers le tutoiement car l'inquiétude prend le dessus}
Tu es aussi blanche que la neige qui commence à tomber ce soir, et tu sembles être en sueur.
On croirait que tu as vu un démon, fit alors Evangeline.

Un démon.
Oui, Seraphina avait bel et bien vu un démon, mais à qui pouvait-elle le dire ?
Avec qui pourrait-elle partager son fardeau ?
Elle ne pouvait pas infliger cela à Evangeline.
\rep Veux-tu qu'on aille au salon de musique ?
Nous serions plus au calme, et on pourrait finaliser les détails de notre prestation, reprit Evangeline.
Ou si tu préfères, on peut rejoindre ma chambre quelques instants pour te permettre d'être au calme d'ici à l'arrivée du Dr Harrington qui devrait, je l'espère, ne point tarder à arriver, conclut-elle.

Non, l'innocence d'Evangeline ne méritait décidément pas d'être entachée par ses égarements.

\fin

\chapitre{Seule}

\noteauteur{Chapitre de flashback quelques années plus tot (environ 4-5 ans). Seraphina ne connait pas encore Florence, et son duo avec Evangeline est recent.
Elle connait deja Beatrice, mais celle-ci est absente. Seraphina ne possede pas encore les partitions de Valentini.}

\debut{W}{itechapel}, Londres, 188?.
\areprendre{Definir la date}
Cela faisait maintenant quelques heures que Mary grelottait dans les rues froides et crasseuses de Whitechapel.
Elle tenta douloureusement de remuer ses doigts engourdis par l'atmosphère glaciale de cette soirée, mais ceux-ci refusèrent de lui répondre.
Même constat du côté de ses pieds.
Son estomac, désespérément vide ces derniers temps, ne se manifestait même plus.
Il fallait à tout prix qu'elle se lève.
Mais le froid, la faiblesse de son maigre corps, les vertiges à chaque pas, les halos blancs masquant son champ de vision, c'était trop pour elle.
Elle n'avait plus la force de faire ne serait-ce qu'un pas.
Et pourtant, il le fallait.
Car, comme si le froid et la faim ne suffisaient pas, cette partie de Whitechapel était le théâtre de nombreuses disparitions ces derniers temps, au moins deux par semaine, en ne comptant que celles qui avaient été signalées.
\noteauteur{Les disparitions frequantes sont due à un pic d'activité de Marcus et son reseau, pour la plupart}
Uniquement des personnes jeunes, orphelins pour la plupart, et qui n'avaient personne pour se soucier de leur absence.
Comme elle.

Comme \areprendre{Repetion de "comme"} si cela pouvait l'aider à reprendre des forces, elle tenta alors de se remémorer à quand remontait son dernier véritable repas.
Si elle ne se trompait pas, c'était il y a plus ou moins deux semaines \areprendre{ajuster la durée pour la coherence avec son état phy et psy}, lorsque Béatrice lui annonça qu'elle devait rentrer quelque temps à Paris pour honorer son défunt père et régler sa succession.
Depuis qu'elles s'étaient croisées par hasard après une représentation au St James’s Hall, à Piccadilly, Béatrice l'invitait régulièrement dans sa maison de Bloomsbury.
Les deux femmes devaient se montrer prudentes, car si Edouard, le mari de Béatrice, avait connaissance de cette relation amicale et de la présence de Mary chez lui, qui aurait pu prédire quelle réaction il serait capable d'avoir, quelles atrocités il serait capable de commettre, que ce soit envers Béatrice ou Mary ?
Elles étaient conscientes du risque, et avaient décidé de l'accepter.
Cela permettait à Mary de manger à sa faim, et cela semblait offrir à Béatrice un peu de sérénité, lui permettant l'espace de quelques heures d'oublier les persécutions que lui faisait subir Edouard.
Mais pour le moment, il allait falloir tenir sans elle, jusqu'à son retour à Londres dans quelques semaines.
Et les maigres restes qu'elle avait pu glaner ici ou là depuis son absence ne suffiraient pas à rester en vie jusqu'à ce moment.
Sans parler du froid.
Et des disparitions.

Mary avait bien pensé à demander de l'aide à Evangeline, mais elle et sa famille passaient l'hiver à la campagne, dans leur manoir près de Bath.
Et de toute manière, qu'aurait-elle bien pu leur dire ?
Ils avaient beau être au courant qu'elle ne venait pas de l'aristocratie et qu'elle avait des origines plus modestes qu'eux, ils n'imaginaient pas à quel point.
Bien qu'ils aient pour le moment l'air de ne pas s'en soucier, Mary redoutait leur réaction s'ils apprenaient qu'elle n'était qu'une simple orpheline de Whitechapel.

Elle releva péniblement sa tête qu'elle maintenait jusqu'alors sur ses genoux pour tenter de conserver un peu de chaleur, et s'aperçut que la neige tombait.
Pas cette jolie neige blanche et douce qui recouvrait délicatement les beaux quartiers de Londres, mais cette neige sale, noire, dont les flocons étaient déjà maculés de charbon avant même de toucher le sol.
\pensee{Comment garder espoir dans ces conditions ?} se demanda-t-elle.
\pensee{Comment garder espoir quand seul un miracle pourrait me permettre de ne serait-ce que survivre jusqu'à demain ?}
Elle décida alors de fredonner un air qu'elle avait écrit il y avait quelques années, le soir où elle fut chassée de l'orphelinat, et qu'elle avait interprété avec Evangeline quelques jours avant son départ, lors de leur dernier concert de la saison, pour essayer de penser à autre chose, ne serait-ce que l'espace de quelques instants.

\separation

La neige ne tombait que depuis une heure au-dessus de Londres, mais déjà, une couche de neige importante pour la saison s'était formée.
Ailleurs dans la ville, ce spectacle aurait pu être celebré pour sa beauté, mais a Whitechapel, même la neige ne pouvait échapper à la corruption et à la noirceur.
Ici, la neige n'était qu'un fléau de plus, dans un quartier qui en compte déjà trop.

Mary était en train d'enlever la neige qui s'accumulait sur ses vêtements, quand soudain, un crissement dans la neige se fit entendre.
Elle s'immobilisa immédiatement, et s'arrêta en plein milieu du refrain qu'elle fredonnait.
Un second craquement lui parvint du bout de la ruelle.
\rep Il\ldots{} il y a quelqu'un ?
lança-t-elle, d'une voix qu'elle voulait franche et assurée, mais qui lui parut terriblement décevante.

Une voix, vraisemblablement masculine, mais étouffé dans un foulard ou un tissu recouvrant la bouche de la personne, lui répondit alors.
\rep Excusez-moi, je ne voulais pas vous faire peur \noteauteur{Le vouvoiement est-il coherent dans ce contexte ?}.
Je vous ai entendu chanter, et je me demandais ce qui pouvait vous pousser à rester dehors par ce temps, si vous aviez besoin d'aide.

L'homme, sans même attendre la réponse de Mary, s'approcha, sans precaution cette fois, et vint s'agenouiller délicatement à côté d'elle.

Tandi qu'il prenait place, Mary tourna la tête vers lui.
L'homme portait un foulard sur le bas de son visage, probablement pour se protéger du froid et de l'atmosphère remplie de crasse.
Il portait également une casquette semblable à celle qu'utilisent bon nombre de ramoneurs.
Mais, bien que rien dans la tenue de l'homme ne soit particulièrement extravagant pour le quartier, Mary, qui avait pu frequenter les hauts cercles de Londres durant l'été, remarqua immédiatement que les tissus composant les vêtements de l'étrange individu était de qualité bien supérieure à ce que l'on pouvait trouver de mieux à Whitechapel.
Baissant alors le regard vers les mains de l'homme, elle aperçut que ce dernier portait des mitaines, masquant la majeure partie de ses mains, et la noirceur de la nuit ne permit pas à Mary de determiner si ces mains semblait avoir connu le dur labeur qu'exigent les quartiers pauvres de Londres pour survivre.
Remontant la tête, elle se rendit compte qu'elle était incapable de voir le visage de l'homme.
Ce dernier tenterait-il de masquer son visage ?
\pensee{Non}, se dit elle, \pensee{même à visage decouvert, l'obsurité m'empecherait de distinguer ses traits.}
\pensee{Et puis, de nombreuses personne qui peuvent se le permettre porte un tissu ou un foulard sur le visage pour proteger leur systeme respiratoire de la saleté ambiante.}

Après quelques minutes sans dire un mot, l'homme fit enfin :
\rep Le petit air que tu fredonnais tout à l'heure, il est très joli !
C'est toi qui l'as composé ?

Mary n'osa pas répondre de vive voix, mais fit un leger hochement de tête.
Cela sembla satisfaire l'homme, qui reprit :
\rep Rassure-toi, je ne te veux pas du mal, tu sais ?
Si j'avais voulu te faire quoi que ce soit, ne crois-tu pas que je me serais deja mit à l'œuvre ?
Je me balade régulièrement dans le quartier pour tenter d'aider ceux que je croise.
Tu es consciente que tu risques de mourir si tu passes la nuit comme ça ?
Si tu ne te fais pas enlever avant.

Bien que l'homme parut sincère, Mary ne sut pas quoi penser de son discours.
En effet, s'il voulait lui faire du mal, pourquoi n'aurait-il pas deja commencé ?
Mais en même temps, pourquoi un homme visiblement riche se souciait-il du sort des habitants de Whitechapel ?

\rep Quel est ton prenom, chère enfant ?
reprit l'homme.

Elle hésita quelques instants, avant de finalement lâcher, d'une voix tremblante :
\rep Mary\ldots{} Et vous ?
\rep Mon nom n'a pas d'importance, Mary.
Dit moi, tu n'as nulle part où aller ?
As-tu seulement mangé à ta faim ces derniers jours ?

Mary ne sut que répondre, mais cette confirmation extérieure de sa solitude la fit éclater en sanglots.
L'homme prit alors sa tête contre lui, comme pour la consoler.
Elle remarque alors que les vêtements de ce dernier dégageaient une forte odeur qui lui était inconnue.
\areprendre{Une phrase qui decrit l'odeur du formol}
\rep C'est tout, calme-toi.
Si tu me fais suffisamment confiance pour me suivre, je peux t'amener jusqu'à une petite auberge tenue par une connaissance qui pourra t'accueillir autant de temps qu'il le faudra.
C'est un peu rustique et austere, mais au moins, tu auras un toit et de quoi manger.
Je t'aurais bien invité chez moi, mais je crains de n'être dans l'incapacité de t'héberger pour le moment.

%elipse, marcus (avec casquette basse et foulard sur le visage) approche (il a initialement l'intention de l'enlever, mais il reconnait Sera (grace a l'air qu'elle fredonne), et se dit qu'elle sera plus utile en vie), l'aide et lui donne quelques pieces.
%Sera se rejouit de pouvoir repondre a la lettre qu'elle a recu d'eva (elle n'avait pas de quoi poster une reponse), mais se demande ce que va lui couter cette aide providentielle.
%Les partitions ne sont pas encore evoquées, Marcus prefere prendre le temps de la mettre sous son emprise pour maximiser les chances qu'elle accepte.
%ni Sera ni le lecteur ne peuvent identifier Marcus. Marcus ne dit et montre pas qu'il a reconnu Sera (Sera ne sait pas que c'est son identité qui l'a sauvée)
%si ça passe mieux en terme de rythme, ce second segment peut être deplacé dans un chapitre a part, a voir en fonction des chapitres du present.

\fin
